\section{Introduction}
\label{chp:intro}

\subsection{The Problems}

Memory is the central resource of modern computing systems.
% TODO phrase next sentence better
Every workloads ultimately reads and writes memory;
every hardware advance in processors, accelerators, and interconnects must be reflected in
how the software stack manages, protects, and translates memory.
Yet, the OS and system software layers that manage memory are among the oldest and most
hardwired components of a modern computing stack.

They were designed for a world of megabytes of DRAM, single-socket servers, and batch workloads.
Today's systems look nothing like that world:
cloud servers carry terabytes of memory,
and virtual machines host DMA-capable devices run with hundreds of gigabits of interconnect bandwidth,
and AI workloads powered by tens of billions of parameters require tens of gigabytes of GPU memory.



This dissertation argues that memory management system---the software machinery responsible for
mapping and protecting virtual memory---is the critical gap between
where hardware capabilities now stand and where application performance can actually go.
The mismatch manifests in three distinct but related problems: 
spanning across CPU memory to I/O memory and device memory.
\begin{enumerate}
\item \textbf{CPU virtual memory translation} has become a major bottleneck for
    memory-intensive workloads on terabyte-scale servers, yet commodity OS kernels
    remain hardwired to a single translation architecture---the radix-tree page table---and
    provide no extensible interface to experiment alternatives designs.
    The difficulties of building OS support has discouraged disruptive architecture research.

\item \textbf{I/O memory protection in virtualized environments} is prohibitively expensive.
    Existing IO memory protection mechanisms offer powerful, so-called strict, safety guarantees; 
    however, they require at least one VM exit per DMA in virtualized environments 
    to invalidate translation entries in IOMMU caches. 
    The VM exits result in 83--95$\%$ application-level throughput degradation ---
    a penalty so severe that production deployments universally disable strict protection, leaving
    them vulnerable to security risks.
    % today forgo protection entirely and accept the associated security risk.

\item \textbf{GPU memory oversubscription for memory-intensive AI workloads}  
    is achieved by each framework implementing its own complex, application-specific offloading.
    CUDA Unified Virtual Memory (UVM) offers a driver-level transparent alternative,
    but is considered too slow for production,
    because the runtime has no visibility into high-level access patterns.
    % too slow, because the runtime has no visibility into high-level access patterns.
\end{enumerate}

\subsection{Thesis Statement}

Modern memory systems need re
% Each problem stems from the same root cause:
% the existing OS or runtime abstraction was designed without the hardware
% or workload it now faces, and neither the interface nor the implementation
% is extensible enough to absorb the mismatch gracefully.
% In each case, the cost is paid in performance, security, or engineering complexity.

% This dissertation addresses each problem with a system that introduces
% a targeted OS/system-level mechanism to close the gap---not by
% bypassing the existing stack, but by making it aware of the hardware
% and workload characteristics it was missing.

\paragraph{EMT: An OS Framework for New Memory Translation Architectures.}
The first problem---hardwired radix-tree translation in Linux---blocks researchers and
hardware designers from experimenting with new translation architectures.
New schemes such as hash-table-based ECPT~\cite{skarlatos2020ecpt} and
flat FPT~\cite{park2022fpt} have been shown to reduce translation overhead
significantly in simulation, but almost no proposals have ever run on a commodity OS.
Without OS experiments, hardware research relies on performance models that assume OS overhead
is constant across translation architectures---an assumption we show to be false.

We present \emph{Extensible Memory Translation (EMT)}, a pragmatic framework
built atop Linux that separates translation-architecture-specific logic from the
architecture-independent memory management core.
EMT provides a stable, hardware-neutral interface through which an MMU driver
can plug in a different page table structure
without touching architecture-independent code.
We port Linux's memory management subsystem onto EMT (EMT-Linux),
showing that extensibility can be achieved with less than 0.5\% overhead on average.
Using EMT, we build full OS support for both ECPT and FPT on Linux,
and evaluate them end-to-end with an emulated MMU in QEMU.
Our experience surfaces correctness and performance challenges
in these hardware designs that are invisible to trace-driven or model-based evaluation,
and we reflect on implications for future hardware/OS co-design.

\paragraph{Virtually Free: Making Strict IO Memory Protection Feasible for Virtualized Environments.}
The second problem concerns IO memory protection in virtualized environments.
Device pass-through delivers near-native I/O performance to virtual machines,
but allowing a guest-controlled device direct access to physical memory poses severe
intra-guest security risks.
Virtualized IOMMUs (vIOMMUs) address this by exposing a virtual IOMMU to the guest,
enabling fine-grained DMA protection.
However, enforcing the highest level of protection---\emph{strict} invalidation,
which purges IOTLB entries immediately after each DMA---requires
a synchronous VM exit per unmap in current hardware-assisted designs.
A VM exit costs 10--20\,$\mu$s; a bare-metal IOTLB invalidation costs 200--300\,ns.
The result is 83--95\% application throughput degradation---a penalty so severe that
production deployments universally disable strict protection.

We present a redesign of strict I/O memory protection that reduces this overhead
to 1--8\% of the unprotected system throughput.
The key insight is that the current Linux IOMMU subsystem serializes all concurrent unmap
requests across cores on a global lock, compounding VM exit latency with lock contention.
Strict security requires only a per-thread ordering
(\emph{Unmap} $\to$ \emph{Invalidate} $\to$ \emph{Free}), not global serialization.
We introduce two mechanisms:
\emph{one-for-many invalidations}, which delegates IOTLB invalidation to a single core and
batches concurrent requests, eliminating redundant VM exits;
and \emph{Deferred Free Page (DFP)}, a callback-based interface that lets drivers delay
physical page reclamation until hardware invalidation is confirmed,
removing synchronous waits without relaxing the security invariant.
Together, these mechanisms achieve performance within 1--8\% of an unprotected system---a
74--147$\times$ improvement over state-of-the-art strict vIOMMU---while maintaining the
same strict security guarantees.

\paragraph{PyTorch-UVM: Practical Unified Virtual Memory for Large Language Model Inference.}
The third problem is GPU memory management for LLM inference.
Flagship model sizes grow faster than GPU memory capacity,
yet the transparent alternative provided by the CUDA driver---Unified Virtual Memory (UVM)---has
been considered impractical for production: prior work reports UVM is more than
5$\times$ slower than application-level offloading.
The gap has two sources.
First, ML frameworks like PyTorch bypass UVM with their own caching allocators,
so managed allocations miss pooling and reuse optimizations.
Second, the driver has no visibility into model execution order,
so page faults and migrations are triggered late and often evict data that is about to be reused.


We present \emph{Pytorch-UVM}, an application-aware UVM substrate
integrated directly into PyTorch.
Pytorch-UVM contributes three mechanisms.
A \emph{UVM-aware caching allocator} proactively returns cached tensors to the driver
when the GPU is oversubscribed, preventing the OOM kills that naive UVM caching suffers.
\emph{Discard-based eviction} marks cached blocks as reclaimable without copying
their data back to host memory, cutting eviction cost.
For workloads where the active working set substantially exceeds GPU memory,
we study \emph{direct CPU-memory access} and
\emph{access-counter-based migration} to break the eviction--re-fault cycle.
Under moderate oversubscription, where the native allocator fails with OOM,
Pytorch-UVM remains 45$\times$--58$\times$ faster than default UVM and
27$\times$--35$\times$ faster than HuggingFace KV-cache offloading,
while matching the native \texttt{cudaMalloc} allocator within 20\% when
memory is not oversubscribed.

\paragraph{Summary of Contributions.}
This dissertation makes the following contributions:

\begin{itemize}
    \item \textbf{EMT} (\emph{EMT: An OS Framework for New Memory Translation Architectures}):
        A framework that gives Linux an extensible, architecture-neutral
        interface for memory translation, enabling OS support for radically different
        page table structures without forking the kernel; plus an open platform for
        evaluating OS kernels on new translation hardware.

    \item \textbf{Virtually Free} (\emph{Virtually Free: Making Strict IO Memory Protection
        Feasible for Virtualized Environments}):
        A redesign of the Linux IOMMU subsystem that replaces
        global cross-core serialization with per-thread safety ordering,
        making strict I/O memory protection feasible in virtualized environments for
        the first time without sacrificing throughput.

    \item \textbf{PyTorch-UVM} (\emph{PyTorch-UVM: Practical Unified Virtual Memory for
        Large Language Model Inference}):
        An application-aware UVM integration for PyTorch that
        propagates model-execution intent to the CUDA memory driver,
        turning UVM into a practical drop-in replacement for bespoke LLM offloading pipelines.

    \item A unifying perspective: across all three systems, the bottleneck is not the
        hardware itself but the OS or runtime abstraction layer that sits between hardware
        capabilities and application demands.
        Each contribution demonstrates that targeted, minimal changes to the right OS interface
        can unlock performance and security gains that no amount of application-level tuning
        can achieve alone.
\end{itemize}

The remainder of this dissertation is organized as follows.
Chapter~\ref{chp:emt} presents EMT and its evaluation on new translation architectures.
Chapter~\ref{chp:viommu} presents the redesigned vIOMMU mechanism and its evaluation
on 400\,Gbps networking workloads.
Chapter~\ref{chp:uvmtorch} presents UVMTorch and its evaluation on LLM inference across
A100, H100, and GH200 hardware.
Chapter~\ref{chp:conclusion} concludes and discusses future directions.
